\documentclass[]{../../../../tex/classes/styledArticle}
\usepackage{../../../../tex/packages/hebrewSupport}

\author{שחר פרץ}
\title{\textit{היסטוריה 43}}
\begin{document}
	\maketitle
	בנוגע להכרזת העצמאות – זה לא היה מובן מאליו. היה ``יותר הגיוני'' לקבל את ההצעה האמריקאית ולא להכריז על המדינה, מאשר להכריז על המדינה. 
	לא סיימנו את פרק ג', אבל שרית סיימה ללמד את פרק ג', בהינתן לחץ הזמנים. נחזור על מבנה הבחינה: 
	\begin{enumerate}[A']
		\item טוטליטאריות ושואה – סיימנו
		\item לאומיות וציונות – סיימנו
		\item בונים מדינה – למדנו על המלחמה במבט על, ולא למדנו על: 
		\begin{itemize}
			\item בעיית הפליטים
			\item הדרוזים
			\item הסכמי שביתת הנשק
		\end{itemize}
		\item סוגיות – לא התחלנו. כולל שלושה נושאים: 
		\begin{itemize}
			\item דה‏־קולוניזציה, וגורל יהודי אצרות האסלם – על זה תינתן משימה לחנוכה
			\item עלייה וקליטה – לא נלמד. ''אני לא אספיק בחיים ללמד עליה וקליטה, זה נושא מאוד מאוד עמוס``
			\item מלחמת ששת הימים/כיפור
		\end{itemize}
	\end{enumerate}
	או במילים אחרות – ''אני לא מלמדת היסטוריה מזמן, אני מלמדת לבגרות``
	
	\section*{תהליך הדה־קולוניזציה}
	נגדיר כמה מושגים: 
	\begin{itemize}
		\item \textbf{קולוניה: }מושבה, חבל ארץ הנשלט ע''י מעצמה אחרת ועל פי רוב מנוצל על–ידה ולצרכים כלכליים וצבאיים. מקור המושג באימפריה הרומית. 
		\item \textbf{קולוניאליזם: }השתלטות של מעצמות אירופה על שטחים מחוץ לאירופה, בשטחי אסיה, אפריקה ואמריקה בסוף המאה ה־19 ותחילת המאה ה־20. 
	\end{itemize}
	המעצמות הקימו בשטחי שליטתן אולוניות, מושבות, ומערכת שלטונית. בחלק מהאימפריות, לקולניות היה עניין תרבותי – לתרבת את האנשים, להעביר להם את הנצרות. בקולוניות באפריקה, הייתה גם תפיסה גזעית. כמו שלמדנו בחט''ב, השיר ''מסע האדם הלבן``. קולוניאליזם זו תפיסה, זו אידאולוגיה. 
	
	\begin{itemize}
		\item \textbf{צרפת} שולטת במרוקו, אלג'יר, טוניס, סוריה ולבנון. במדינות האלו מדברים צרפתית עד היום. לימוד הצרפתית היה חובה על כל מי שחי במדינה. התפיסה היא מעורבות תרבותית, והמעורבות היא מעורבות ישירה. היו ממש שכונות של צרפתים, שהגיעו מצרפת לחיות במדינות האלו, ללא קשר לתפקידם בשלטון המקומי. התפיסה הייתה שהקולניות הן חלק אינטגרלי מצרפת. צרפת שלטו לא רק בעזרת מנדטים. לא הייתה הכרה באליטות פנימיות כמנגנוני שלטון, והם שלטו באופן ישיר באוכלוסיה. 
		\item \textbf{בריטניה: }מצריים, א''י, ירדן, עיראק. שם רק מי שהיה צריך לדעת אנגלית, ידע אנגלית. לחלק מהפקידים היו משפחות, אבל ככלל, בריטניה לא הביאה אוכלוסיה אזרחית. מבחינת בריטניה, היא נמצאת בקולוניות בשביל האינטרסים שלה בלבד, ולא היה את החלק של ''השליחות התרבותית`` הצרפתית. בריטניה שולטת בעיקר בשיטת המנדטים, כלומר הם לא ניהלו ישירות את האוכלוסיה המקומית – הייתה ''הנהגת היישוב היהודי``, והיא התנהלה מול שלטון המנדט הבריטי והחוק הבריטי. בכל מה שקשור לעניינים הפנימיים של היישוב היהודי, הייתה אוטונומיה. בריטניה פתחה שלוש שיטות משטר שונות. 
		\item \textbf{איטליה: }לוב, אתיופיה. 
	\end{itemize}
	
	נגדיר עתה דה־קולוניזציה. \textbf{דה־קולוניזציה} הוא \textit{תהליך שלילת הקולוניאליזם}. התהליך החל עוד בין שתי מלחמות העולם, אך שיאו מתרחב לאחר מלחמת העולם השנייה, ובשנות ה־50-60 של המאה ה־0. הוא התרחש בכל איזור שבו שלטו המעצמות הקולוניאליות. [זו הגדרה שצריך לזכור: אם שוכחים את אחד מהרכיבים, ירד ניקוד]. 
	
	יש מספר גורמים לתהליך הדה־קולוניזציה. ניתן לחלק את הגורמים השונים לשלושה נושאים: 
	\begin{itemize}
		\item גורמים אידאולוגים: 
		\begin{itemize}
			\item \textit{התפתחות רעיון הלאומיות:} השלטון הזר מגיע עם הרעיונות והאידאולוגיה שלו, במיוחד עם הצרפתים. הם מחלחלים באוכלוסיה. חלק מהמקומיים, נלחמים לצד המעצמות. ואז, בדיוק כמו הלאומיות באירופה (אבל בדיליי), נוצרות תנועות לאומיות, ואז מאבק (במקומות האלה, הרבה פעמים אלים מאוד) במעצמות. 
			\item \textit{התנגדות אידאולוגית בשתי מעצמות העל}
			\item \textit{דעת הקהל בעולם:} אחרי מלחמת העולם הראשונה, ובמיוחד אחרי מלחמת העולם השנייה, יש התנגדות לשליטה של עמים על עמים. מבחינה רעיונית, האו''ם קם על הבסיס הזה, של שחרור וריבונות של עמים. 
		\end{itemize}
		\item גורמים כלכליים: 
		\begin{itemize}
			\item \textit{המשבר הכלכלי אחרי מלחמת העולם השנייה: }יש משבר כלכלי. בריטניה וצרפת (שנכבשה) נמצאות במצב כלכלי מאוד קשה. הן יורדות מגדולתן, ולהחזיק קולוניות עולה \$\$\$. המנדט בארץ ישראל איננו קולוניה, אבל בדיוק כמו שעלה כסף להחזיק אותו – עלה כסף להחזיק קולוניות. 
			\item \textit{ירידה ברווחיות המושבות: }הגלובליזציה הולכת ומפתחת, וכבר יש פחות משאבים אקסלוסיביים לאיזורים ספציפיים. 
		\end{itemize}
		\item גורמים פוליטיים: 
		\begin{itemize}
			\item \textit{שינוי יחסי הכוחות בעולם והשפעת המלחמה הקרה:} האינטרס של מעצמות העל הוא שלא יהיו אימפריות שולטות אלא מדינות. 
			\item מאבק התנועות הלאומיות
			\item תפקידו של האו''ם
			\item השלכות מלחמות העולם
		\end{itemize}
	\end{itemize}
	
	אז למה המעצמות לא שחררו ככ מהר? יש להם אינטרסים, גם כלכליים, אבל גם אסטרטגיים. 
	
	מאפייניו הכלליים של התהליך: 
	\begin{itemize}
		\item עמי המזרח התיכון הקדימו את עמי אפריקה, כי הבריטים שולטים בעיקר במזרח התיכון, והעמים תחת שליטת בריטניה מקבלים יותר עצמאות, מכורח אופן שליטת הבריטים, שחשובים להם האינטרסים שלהם גרידא, ולא השליחות התרבותית. לדוגמה, ברגע שבריטניה עושה הסכם על מצרים בנוגע לתעלת סואץ והשיגה שליטה מלאה עליה, היא פשוט שחררה. 
		\item מאבק לאומי קלאסי – גיבוש הרעיון הלאומי ע''י אליטות משכילות, הקמת תנועות, והסתייעות בגורמי חוץ. 
		\item המעצמות ניסו לשמור על האינטרסים שלהם, תוך חתימת ההסכמים על התנועות הלאומיות (לדוגמה סואץ). 
		\item כמו שדיברנו, יש הבדלים בין העמים שקמו תחת השלטון הבריטי ותחת השלטון הצרפתי. 
		\item התהליך הוא תהליך הדרגתי, כמו כל מאבק לאומי. 
	\end{itemize}
	
	כתוצאה מהתהליך הזה השתנתה כליל המפה. הדבר שהכי בולט בה, הוא הקווים הישרים. אין חיה כזאת בטבע, ''קו ישר במפה``. אך, המעצמות חלקו בינהן פשוט עם קווים ישרים, בלי עניין במי ומה חיי שם. חלק גדול מהסכסוכים שאנו מכירים בעולם נובעים מהחלוקה הזאת. ואכן, בערך כל קו ישר במפה החדשה הוא תוצאה של מעצמות שונות שיצאו בזמנים שונים. 
	
	``שאלות. אין? סיימנו את הנושא``. 
	
	הנושא הזה, וגורל יהודי ארצות האסלם, לעיתים קרובות מלווים בשאלות ממש פשוטות. זה נושא ראש בקיר. עכשיו מבחנים. עתה נעבור קצת על המבחן. 
	
	\subsection*{המבחן}
	בפרק ראשון שאלה 1 – הייתה שאלה על היחס ליהודים בגרמניה ובפולין. סעיף א' התייחס לחוקי נירנברג עם שני קטעי מקור ושאלה שמדברת על הערכת המקורות. סעיף ב' דיבר על ההחרפה ביחס ליהודי פולין. מי שענה על השאלה הזו, לרוב ענה יפה על סעיף א'. סעיף ב' לכאורה יותר פשוט, אבל כאן עשו טעויות – במשפט ''בין השנים 39-41 לבין 42-45 הייתה החרפה ביחס יהודי פולין``, וביקשו לתת דוגמאות לשני אירועים. בקשו להסביר את המשפט, אבל היו אנשים שלא הסבירו, כלומר פשוט נתנו דוגמאות לשני אירועים, אבל לא הסבירו את המשפט ישירות. הפואנטה היא החרפה בתקופת הפתרון הסופי. לכתוב ההכנסה לגטאות היא החרפה, זה לא קשור לשאלה. שרית קיבלה מקורות של בורות ירי, אבל זה לא מדויק – היו בורות החלו בברית המועצות, ולא בפולין. יש בורות ירי בפולין, ולכן שרית קיבלה את זה חלקית, אבל ''אל תעשו את זה בבגרות, אני לא בטוחה שיקבלו את זה``. 
	
	שאלה 2 – לא הרבה ענו עליה. בסעיף ב' שלה, היה צריך להסביר את אחד המהלכים במלחמת העולם השנייה והשפעתו על גורל היהודים. יש כאלו שכתבו אל־עלמיין. זהו קרב שמחולק לשניים, רק החלק השני ומבצע לפיד רלוונטי. היה אפשר לכתוב גם על כיבוש צרפת. בנוגע לשאלה של האם שואת יהודי צפון אפריקה היא חלק מהשואה, אל תצטטו את ועידת ואנזה – יש כאן שני דברים, ולא אחד. 
	
	שאלה 3 – דרך ההתנגדות עד 42, והדגמה באמצעות שתי פעולות. עד תחילת הפתרון הסופי, הדרך היא כיבוש החיים. לא תנועות נוער. תנועות נוער זו דוגמה לממש את כיבוש החיים. 
	
	בפרק השני, לאומיות וציונית, הייתה שאלה מספר 4 – בשאלה מספר 4 סעיף ב', הייתה קצת הרבה בעיה. ''הציגו נקודות שוני אחת ונקודת דמיון אחת בין התנועה הלאומית עליה למדתם לתנועות לאומיות כלליות``. זו שאלה מסוג השוואה. יש בה השוואה. 
	''אני שמחה שזה קרה לכם, כי זו דוגמה לאיך לא בוחרים, למה לא לבחור. לא היה ילד אחד שענה על שאלה 4 פה וקיבל את מלא הנקודות``. בשאלת השוואה יש \textit{דומה} ו\textit{שונה}, כפול 2 תנועות לאומיות, כלומר 4 דברים לכתוב עליהם. ואז ירד חצי מהניקוד. 
	
	שאלה 5 – סבבה. שאלה 6 – סבבה. 
	
	בפרק השלישי של בונים מדינה, שאלה 7 – הסיבות להקמת תנועת המרי. שאלו לפי דעתכם, מי מבין הצדדים יותר אחראי לעזיבתם של הבריטי. הכוונה בניסוח פה היא לצורת המאבק – צבאי, העפלה, התיישבות. מי שכתב רצוף/צמוד ונימק את זה לפי צבאי וכו', זה בסדר. מי שלא, פחות. עוד משהו – קטע מקור חזותי, לעיתים קרובות קשה לכם להגיד עליו יותר מדי. אלא אם מבקשים מכם מקור חזותי, תנסו לבדוק אם יש מקור אחר, כי תמיד צריך שני רכיבים בהסבר. לגבי מאבק צבאי, אם אין הגבלה כרונולוגית, אפשר להשתמש בפעילות של תנועת המרי (פה לא הייתה הגבלה). 
	
	שאלה 8 – העברת א''י לאום. שכתב את אקסודוס, והסביר את זה בהסבר של ההשפעה על ועדת אונסק''ו, זה בסדר, כי הניסוח הוא מי אחראי לעזיבתם של הבריטי, ולא למי אחראי להעברת שאלת א''י לאום. אבל תשימו לב לדקויות האלה. 
	
	''חלק מכם קיבלו ציונים לא משהו. יש מתכונת. חלק קיבלו ציונים יפים. וחלק יפים מאוד. ``
	
	
\end{document}